\begin{abstract}
This volume unifies the five papers of the Songlines research
programme into one document. The programme asks how agents with
symbolic semantic memory complete navigation tasks --- and what
changes when memory becomes collective. \textbf{Part I} introduces the
Q/R/M/C stage decomposition (Query $\to$ Retrieval $\to$
Materialization $\to$ Completion): one logging contract that
localises failures which success rates cannot see, portable across
five system classes from symbolic single agents to LLM stacks; its
sharpest single-agent finding is that candidate \emph{generation},
not ranking, is the binding constraint ($91\%$ empty retrievals at
$96.9\%$ ranking precision). \textbf{Part II} makes ``navigating by
someone else's memory'' a measured event: every commitment is
annotated with the foreign share $\varphi$ of its evidence, raw
foreign locks almost never complete directly yet pay off as
transport, a closed-form distance law predicts intention horizons
tick-exactly, and a two-line coordination contract (Warp Drive) turns
the fastest sharing regime into the best one. \textbf{Part III}
identifies the transferable unit: not the pin (a coordinate) but the
song --- routes with edge provenance that inherit the witness's
completion and safety and rupture at a predicted edge; place identity
recovered from meaning alone (failing closed where coordinates fail
open); and the anatomy result that nodes carry identity while edges
carry transport, at $6.4\%$ of snapshot bits. \textbf{Part IV} gives
collective memory its mechanics --- merge, trust, and staleness as
three explicit rules occupying a new regime on the
$M{\times}C$ plane --- and its mathematics: a categorical frame in
which the alignment morphisms are constructed rather than assumed and
the trust$\times$cadence phase diagram decides whether the colimit
spans the collective. \textbf{Part V} turns to the
\emph{ontogenesis} of memory: formation as a two-axis decision
(counterfactual marginal utility $\times$ simplicity of analogy) with
an exception operation that protects against false generalisation ---
measured deterministically, learned by a contextual bandit that
improves on its own designer, and evolved to a finite couplet budget.
Part V then composes all mechanisms into one runtime and closes a
reviewer's acceptance test: a runtime utility estimator that matches
the exact counterfactual (Spearman $0.989$) so the controller needs
no oracle; receiver-validated admission that makes social exchange
net-positive at horizon where unregulated exchange is harmful; an
equal-budget comparison where raw history does not catch up even at
$4\times$ the memory; an adversarial-provenance test; safety gates
driving fail-open below $1\%$ on both a grid and a continuous
substrate; and a full wire-cost accounting. The parts were written as
stand-alone submissions; cross-references between them appear here as
parts of one volume.
\end{abstract}

\section*{How to read this volume}

\paragraph{One sentence per part.} Part I builds the instrument
(stage-decomposed measurement); Part II uses it to price foreign
evidence (provenance and its law); Part III finds the right unit of
transfer (the song: landmarks $+$ beats); Part IV supplies the
collective mechanics and the mathematics of ``meaning the same
thing''; Part V asks how the memory that Parts I--IV coordinate comes
to exist at all, and makes its formation rule learnable and evolvable.

\paragraph{Lineage.} Parts I and IV grew out of a single earlier
manuscript (\emph{Q/R/M/C Framework $+$ Minimal Collective Semantic
Memory}, kept in the series archive as \texttt{symbolic\_memory});
every experiment of that ancestor appears here, refined, in Parts I
and IV. Parts II and III are companions built on Part I's
instrumentation. Where a part's text says ``the companion paper,'' it
refers to another part of this volume.

\paragraph{Measurement discipline.} Every claim in every part follows
the same protocol: predictions registered to disk before execution,
failures against registration reported with their mechanisms, and the
strongest results stated as exact-match validations of closed-form
predictions (the distance law's $6/6$ hold-out, the rupture law's
$110/110$, the $SE(2)$ recovery's $80/80$, the vocabulary ablation's
$0$ fail-open in $252$ cells).

\paragraph{Code and data.} All experiments live in one repository
(\texttt{experiments/}, \texttt{songline\_drive/}); each part's
reproducibility appendix lists its exact commands. Raw run artifacts
are archived under \texttt{tmp/} with per-experiment
\texttt{RESULTS.md} files.
